% p3-01-introduction

\section{P3 §1. Introduction}

1. Introduction

      Figure (fig-p3-open). Opening triptych --- Map / Ladder of Evidence L1..L5 / Falsifier Balance.
                       The three-paper programme charted as a compass journey.

  Two recurrent themes in the mathematical analysis of numerical constants are (i) the
  search for compact symbolic representations and (ii) the characterisation of
  approximation hierarchies that interpolate between exact values and economical closed
  forms. The present paper addresses both themes simultaneously through a synthesis of
  two formalisms developed in companion Paper 1 and companion Paper 2 of this trilogy,
  referred to collectively as the GOLDEN BRIDGE (Trinity-Pellis Bridge) program.

  The Pellis expansion (developed in companion Paper 1 of this trilogy) generalises the
  classical Zeckendorf representation and continued-fraction theory by decomposing a
  target real number x as a hierarchical series in powers of $\varphi$-k, where $\varphi$ = (1 + 5)/2 is
  the golden ratio, with coefficients that may themselves encode interference-like
  cancellations at successive levels of the hierarchy. Zeckendorf's theorem --- that every
  positive integer has a unique representation as a sum of non-consecutive Fibonacci
  numbers --- serves as the discrete prototype; the Pellis expansion is its
  continuous-variable and hierarchically-extended descendant. The expansion is formally
  additive and converges for all positive reals, providing arbitrarily fine approximations at
  the cost of increasing hierarchical depth.

  The Trinity monomial system (developed in companion Paper 2 of this trilogy) organises
  real numbers as products n $\cdot$ 3k $\cdot$ $\varphi$p $\cdot$ pim $\cdot$ eq with (n, k, p, m, q) $\in$ Z5 (with n $\geq$ 1).
  These monomials form a multiplicative abelian monoid under componentwise integer
  addition of exponents and serve as coarse-grained landmarks in the real line. The
  adjective "Trinity" refers to the three transcendental generators $\varphi$, pi, e, together with
  the rational-scaling prefactors n $\cdot$ 3k. The integer 3 carries additional algebraic
  significance via the anchor identity $\varphi$2 + $\varphi$-2 = 3, which roots the Trinity monomial
  system in the Lucas ring Z[$\varphi$] (Section 2.5).

  The unification problem addressed here is: given a real number x with a known Pellis
  expansion to depth D, how does one identify the optimal Trinity monomial approximant,
  quantify the residual, and characterise the structural relationship between the two
  representations? Conversely, given a Trinity monomial, what does its Pellis expansion
  reveal about its fine-scale arithmetic structure?

  The paper proceeds as follows. Section 2 establishes notation, fixes the representation
  hierarchy, and introduces the Lucas ring as the discrete state space. Section 3 develops
  the multiplicative grammar of Trinity monomials. Section 4 presents the Pellis
  expansion with full convergence analysis. Section 5 defines the projection map PiL and
  the residual operator. Section 6 develops the Koopman/transfer-operator formalism.
  Section 7 interprets representation quality via the MDL principle. Section 8 states the
  symbolic phase transition conjecture. Section 9 describes the testable computational
  program. Section 10 presents worked examples, including a careful disambiguation of
  two AlphaPhi objects. Section 11 discusses limitations. Sections 12--13 conclude and
  map the PhD program elements to the Paper 3 formalism. Appendices collect formal
  definitions, the reference implementation, notation, the falsification ledger, and the Coq
  provenance map.

  1.1 Relationship to Prior Work
  Zeckendorf's theorem, originally proved by Edouard Zeckendorf and generalised by
  Hoggatt (1972), establishes that every positive integer has a unique representation as a
  sum of non-consecutive Fibonacci numbers; this is the discrete analogue of the Pellis
  expansion. The theory of $\beta$-expansions provides the continuous-real-number
  generalisation in non-integer base beta = $\varphi$: as Rényi (1957) established in Acta
  Mathematica Academiae Scientiarum Hungarica (vol. 8, pp. 477--493), these expansions
  have well-defined ergodic properties; Parry (1960), in Acta Mathematica Academiae
  Scientiarum Hungarica (vol. 11, pp. 401--416), subsequently characterised which bases
  yield finite or periodic digit sequences. The golden ratio $\varphi$ is itself a Parry number (its
  beta-expansion of 1 is periodic), and the Pellis expansion generalises the Rényi/Parry
  setting by allowing non-binary coefficients and hierarchical interference corrections.
  Boyd's work on beta expansions, building on Parry and using Pisot/Salem number
  theory, gives the backdrop for understanding when such expansions are well-behaved.

  Transfer-operator methods in symbolic dynamics originate with Ruelle's
  thermodynamic formalism (Cambridge Mathematical Library, 2nd ed.), which
  developed the mathematical structure of equilibrium statistical mechanics for
  dynamical systems. Baladi's Positive Transfer Operators and Decay of Correlations
  (World Scientific, Advanced Series in Nonlinear Dynamics, vol. 16, 2000) provides the
  functional-analytic machinery --- particularly spectral gap theorems on appropriate
  Banach spaces --- that the spectral gap conjecture of Section 6 would require to be
  established rigorously.

  The MDL principle was introduced by Rissanen (1978) in Automatica and
  comprehensively treated by Barron, Rissanen, and Yu (1998) in IEEE Transactions on
  Information Theory (vol. 44, no. 6, pp. 2743--2760). The connection between MDL and
  Kolmogorov/algorithmic complexity draws on Chaitin's Algorithmic Information Theory
  (Cambridge University Press, 1987); the full MDL treatment by Grünwald (2007, MIT

  Press) provides the modern framework.

  The connection between the golden ratio and quasicrystalline structure --- motivated by
  Shechtman, Blech, Gratias, and Cahn's discovery (Physical Review Letters, vol. 53, no.
  20, pp. 1951--1954, November 1984) of a metallic phase with long-range icosahedral
  orientational order and no translational symmetry --- provides a physically-motivated
  reason for why $\varphi$-based symbolic representations are mathematically natural. In
  Penrose tilings and quasicrystal diffraction patterns, the ratio of interatomic distances
  is always related to $\varphi$, and the Fibonacci sequence governs the ordering of the
  aperiodic structure. The Trinity framework inherits and algebraises this structural role
  of $\varphi$.

  The connection to the Trinity $S^{3}$AI / Flos Aureus PhD program is new to Paper 3. The
  42-entry precision catalogue provides a concrete target set for the computational
  program of Section 9, but it is not treated as validation until the MDL tests are
  executed. The Coq citation map records the formal proof status of algebraic identities
  used here. Exact global theorem counts are intentionally not reproduced in the paper,
  because those counts are regenerated as the proof base evolves. Admitted obligations
  are cited conservatively as "proof-sketch" status throughout this paper.

  1.2 Scope and Honest Boundary
  Per the falsification protocol of the PhD program's Appendix B and Chapters 18 and 52
  (Limitations), all claims in this paper are bounded as follows:

  - Claims marked (Theorem) have proofs in the body or refer to published results.
  - Claims marked (Conjecture) are open; they are stated for their structural value and to
  generate a testable computational program.
  - Claims marked (Working hypothesis) depend on unproven number-theoretic
  independence results (see Section 11.1).
  - Claims marked (Proof-sketch) correspond to identities whose Coq formalisation
  carries one or more Admitted lemmas.
  - No physical constants are asserted to be structurally determined by the GOLDEN BRIDGE (Trinity-Pellis Bridge)
  framework. The framework is a compression language, not a physical theory.

\subsection{References}

- Vasilev, D. \textit{Flos Aureus Monograph: Trinity Constants}. DOI \href{https://zenodo.org/doi/10.5281/zenodo.19227877}{10.5281/zenodo.19227877}.
- CODATA 2018 recommended values: NIST \href{https://physics.nist.gov/cuu/Constants/}{https://physics.nist.gov/cuu/Constants/}.
- Heyrovska, R. Golden-angle FSC publications --- Heyrovský Institute of Physical Chemistry.
